We are now ready to state Green	extquotesingle s formula.

\begin{Theorem}
\label{thm:elliptic-green-formula-adjoint-boundary-operators}
Let $L(\bm{x},D)$ be an elliptic operator of order $2m$ on $\overline\Omega$ and let $B_j(\bm{x},D)$, $j=1,\ldots,m$, be a set of normal boundary operators. Let $S_j(\bm{x},D)$, $j=1,\ldots,m$, be a set of boundary operators which complements the $B_j$ to form a Dirichlet system. Then there exist boundary operators $C_j(\bm{x},D)$, $T_j(\bm{x},D)$, $j=1,\ldots,m$, with the following properties:\begin{enumerate}\item $\operatorname{ord}C_j=2m-1-\operatorname{ord}S_j$, $\operatorname{ord}T_j=2m-1-\operatorname{ord}B_j$. ($\operatorname{ord}$ stands for the order of the operator.)\item The $C_j$ and $T_j$ form a Dirichlet system.\item For every $u,v\in H^{2m}(\Omega)$, we have
\begin{equation}\label{eq:elliptic-green-formula-adjoint-boundary-pairing}
\int_\Omega(Lu)v-u(L^*v)\,\mathrm{d}\bm{x}=\sum_{j=1}^m\int_{\partial\Omega}(S_ju)(C_jv)-(B_ju)(T_jv)\,\mathrm{d}S.
\end{equation}
Here $L^*$ is the formal adjoint of $L$; see Definition~\ref{def:distribution-convolution-formal}.\item If the $B_j$ satisfy the complementing condition for $L$, the $C_j$ satisfy the complementing condition for $L^*$.\end{enumerate}
\end{Theorem}

\begin{Proof}
Integration by parts yields a formula of the form
\begin{equation}\label{eq:green-formula-raw-boundary-derivative-sum}
\int_\Omega(Lu)v-u(L^*v)\,\mathrm{d}\bm{x}=\sum_{\alpha,\beta}\int_{\partial\Omega}a_{\alpha\beta}(\bm{x})D^\alpha uD^\beta v\,\mathrm{d}S,
\end{equation}
where the sum extends over $\alpha$ and $\beta$ with $|\alpha|+|\beta|\leq2m-1$. We next integrate by parts on $\partial\Omega$ and move all tangential derivatives from $u$ to $v$, so that only purely normal derivatives of $u$ are left (carrying out this step requires a partition of unity and local coordinate charts on $\partial\Omega$). This leads to a formula of the form
\begin{equation}\label{eq:green-formula-pure-normal-boundary-reduction}
\int_\Omega(Lu)v-u(L^*v)\,\mathrm{d}\bm{x}=\sum_{j=0}^{2m-1}\int_{\partial\Omega}a_j(\bm{x})\frac{\partial^ju}{\partial n^j}E_j(\bm{x},D)v\,\mathrm{d}S,
\end{equation}
where $E_j$ is a differential operator of order $2m-j-1$. If $L$ is elliptic (or, even more generally, if $\partial\Omega$ is noncharacteristic), then $E_j$ contains a term proportional to $\dfrac{\partial^{2m-j-1}v}{\partial n^{2m-j-1}}$ with a nonzero coefficient; in other words, the $E_j$ form a Dirichlet system. We next use Lemma~\ref{lem:dirichlet-system-tangential-triangular-reduction} to find
\begin{equation}\label{eq:green-formula-normal-derivative-dirichlet-system-expansion}
\frac{\partial^ju}{\partial n^j}=\sum_{k=1}^m\Psi_{jk}(\bm{x},D)B_k(\bm{x},D)u+\sum_{k=1}^m\Psi_{jk}^{\prime}(\bm{x},D)S_k(\bm{x},D)u.
\end{equation}
We substitute this into \eqref{eq:green-formula-pure-normal-boundary-reduction} and then integrate by parts on $\partial\Omega$ to move the tangential differential operators $\Psi_{jk}$ and $\Psi_{jk}^{\prime}$ from $u$ to $v$. This yields \eqref{eq:elliptic-green-formula-adjoint-boundary-pairing}.

To verify the complementing condition, let $\Omega$ be the half-space $\{x_n>0\}$ and let $L$ have constant coefficients. Consider solutions of the form $\exp(i\bm{\xi}\cdot\bm{x})v(x_n)$, where $\xi_n=0$ and $v(x_n)\to0$ as $x_n\to\infty$. Green\textquotesingle s formula \eqref{eq:elliptic-green-formula-adjoint-boundary-pairing} holds with $\Omega$ replaced by $\Sigma\times\mathbb{R}^+$, where $\Sigma$ is a parallelepiped in $\mathbb{R}^{n-1}$ corresponding to one period. Moreover, $L$ now becomes an ordinary differential operator. For such operators, a Fredholm alternative holds, i.e., the initial-value problem
\begin{equation}\label{eq:green-complementing-homogeneous-normal-ode}
L(\bm{x},D)u=0,\ B_j(\bm{x},D)u=0\quad\text{for }x_n=0
\end{equation}
has only the trivial solution if and only if the problem
\begin{equation}\label{eq:green-complementing-inhomogeneous-normal-ode}
L(\bm{x},D)u=f,\ B_j(\bm{x},D)u=g_j\quad\text{for }x_n=0
\end{equation}
is solvable for all $f$ and $g_j$. By Green\textquotesingle s formula, the latter condition implies that the initial-value problem
\begin{equation}\label{eq:green-complementing-adjoint-homogeneous-normal-ode}
L^*(\bm{x},D)v=0,\ C_j(\bm{x},D)v=0\quad\text{for }x_n=0
\end{equation}
has only the trivial solution, i.e., that the $C_j$ satisfy the complementing condition for $L^*$. Note that if $v$ were a nontrivial solution of \eqref{eq:green-complementing-adjoint-homogeneous-normal-ode}, then $f$ and $g_j$ in \eqref{eq:green-complementing-inhomogeneous-normal-ode} would have to satisfy
\begin{equation}\label{eq:green-adjoint-mode-data-compatibility}
\int_\Omega fv\,\mathrm{d}\bm{x}=-\sum_{j=1}^m\int_{\partial\Omega}g_jT_jv\,\mathrm{d}\bm{x}.
\end{equation}
This completes the proof.
\end{Proof}
